\chapter{Conclusion}\label{ch:conclusion}

End-to-end encrypted messaging depends on authentic public keys, yet manual comparison is rarely
completed~\cite{vaziripour2017}. Moreover, the deployed transparency systems surveyed in
\cref{ch:background} do not describe a conversation-specific channel that warns both peers of an
individual key substitution. This thesis addresses that gap through DUET, which automatically
checks for key substitution and provides an in-band peer-warning channel for honestly established
and monitored conversations without changing the lookup-proof format.

DUET instantiates ACKA's abstract log component $\Log$ as a verifiable key directory designed to
expose the malicious-operator behaviours specified in the threat model (\cref{ch:threat}). It adds a
dual-index warning scheme for in-band peer
warnings and a windowed consistency proof for offline catch-up. DUET also retains ACKA's CEA
construction in continuous mode, using simulated ratchet inputs in the prototype. A separate auditor
process verifies directory transitions. Detecting split-view equivocation in deployment requires
comparison of signed heads obtained from independent vantage points. The \texttt{duet} Rust crate
implements the directory server, proof construction, and auditor process. A verifier integrated into Signal Desktop runs the client-verification path
(\cref{ch:impl}).

\Cref{ch:eval} shows that detection uses additional label-bound lookups without enlarging an
individual lookup proof. At a matched directory size, DUET's measured lookup proof and the derived CONIKS quantity are in the
same sub-kilobyte size class when their signed heads are excluded. Median full verification
remains below \SI{160}{\micro\second} through $10^7$ labels. Windowed catch-up uses one consistency
proof for the complete missed window. These measurements cover local components and exclude
end-to-end monitor latency. The adversarial suite produced the expected outcome in all eight specified cases. The STRIDE analysis identified one privacy residual and three unimplemented
availability controls.

These results support the four contributions introduced in \cref{sec:intro-contributions}.
Conceptually, the dual-index scheme gives both parties access to a shared warning location while
retaining party-specific lookup labels. Cryptographically, DUET applies an RFC~9162
distant-head consistency proof to offline monitoring. The DUET directory server and Signal Desktop verifier show that client-side verification is feasible. Empirically, the benchmark suite quantifies the component costs, and the adversarial suite tests the
controls across the attack scenarios.

Production integration still requires binding the per-session labels to Signal's live conversation
secret and defining how failed checks affect session establishment (\cref{sec:disc-future}). Under
the assumptions of \cref{sec:threat-assumptions}, the prototype shows that these automatic checks can
run inside a messaging client and, in the demonstrated dual-index warning path, notify both parties of
a detected substitution.
